\documentclass[12pt,a4paper]{article}

\usepackage[margin=2.5cm]{geometry}
\usepackage{amsmath,amssymb,amsthm}
\usepackage{setspace}
\usepackage{natbib}
\usepackage{graphicx}
\usepackage{hyperref}
\usepackage{booktabs}
\usepackage{enumitem}
\usepackage{tcolorbox}
\usepackage{mdframed}
\usepackage{xcolor}

% Define missing colors
\definecolor{darkgreen}{rgb}{0.0, 0.5, 0.0}

% Extend appendix counter beyond Z using double letters
\makeatletter
\newcommand{\AlphExtended}[1]{%
  \ifcase\value{#1}\or A\or B\or C\or D\or E\or F\or G\or H\or I\or J\or K\or L\or M%
  \or N\or O\or P\or Q\or R\or S\or T\or U\or V\or W\or X\or Y\or Z%
  \or AA\or AB\or AC\or AD\or AE\or AF\or AG\or AH\or AI\or AJ\or AK\or AL\or AM%
  \or AN\or AO\or AP\or AQ\or AR\or AS\or AT\or AU\or AV\or AW\or AX\or AY\or AZ%
  \else\@ctrerr\fi
}
\makeatother

% Custom environments for appendices
\newtcolorbox{predictionbox}[1][]{colback=blue!5,colframe=blue!50!black,title=#1}
\newenvironment{prediction}{\begin{mdframed}[backgroundcolor=blue!5]}{\end{mdframed}}
\newenvironment{hypothesis}{\begin{mdframed}[backgroundcolor=green!5]}{\end{mdframed}}
\newenvironment{falsificationbox}{\begin{mdframed}[backgroundcolor=red!5]}{\end{mdframed}}
\newenvironment{testbox}{\begin{mdframed}[backgroundcolor=yellow!5]}{\end{mdframed}}
\newenvironment{keyinsight}{\begin{mdframed}[backgroundcolor=blue!10]}{\end{mdframed}}
\newenvironment{noveltybox}{\begin{mdframed}[backgroundcolor=orange!10]}{\end{mdframed}}

% Unicode checkmark support
\usepackage{newunicodechar}
\newunicodechar{✓}{\checkmark}

\setstretch{1.2}
\bibpunct{(}{)}{;}{a}{}{,}

% Theorem environments
\newtheorem{theorem}{Theorem}[section]
\newtheorem{proposition}[theorem]{Proposition}
\newtheorem{lemma}[theorem]{Lemma}
\newtheorem{corollary}[theorem]{Corollary}
\newtheorem{definition}[theorem]{Definition}
\newtheorem{example}[theorem]{Example}
\newtheorem{remark}[theorem]{Remark}
\newtheorem{axiom}[theorem]{Axiom}

\title{Complementarity and Context:\\
A Unified Framework for Economic Rationality}

\author{Gerhard Fehr\\
FehrAdvice \& Partners AG, Zurich\\
Beatrix Research Group}

\date{January 2026 -- Version 51}

\begin{document}

\maketitle

\begin{abstract}
Economics remains fragmented: classical theory explains market efficiency but not social preferences; behavioral economics explains social preferences but not institutional change; institutional economics explains institutional change but not financial crises. This paper provides a unified framework based on \emph{complementarity}---the interdependence of choices across agents, domains, and time---and \emph{context}---the technological, institutional, normative, and cultural environment that shapes payoff structures.

The framework is proposed as a \emph{metatheory}: rather than replacing established theories, it aims to subsume them as limiting cases and specify \emph{when} each applies. Arrow-Debreu emerges when complementarity is zero (see Appendix~A for formal derivation) ($C_{ij} = 0$). Behavioral economics emerges when context is fixed ($\Psi = \bar{\Psi}$). Institutional economics emerges when complementarity is simplified. The framework's value-added lies in predicting which theoretical tradition applies to which situation---enabling conditional rather than unconditional predictions.

We introduce four core concepts: (1) the complementarity matrix $C$, capturing how choices reinforce or undermine each other; (2) the context vector $\Psi \in \mathbb{R}^8$, decomposed into eight primitive dimensions; (3) the reference structure $C^*(\Psi)$, derived as the self-consistent fixed point toward which systems gravitate; and (4) the coherence index $K$, measuring alignment between actual and reference structures.

\textbf{Key methodological contribution:} We operationalize context through eight independently measurable dimensions:
\begin{itemize}
    \item $\Psi_I$ (Institutional): Rules, enforcement, property rights
    \item $\Psi_S$ (Social): Trust, norms, social capital
    \item $\Psi_C$ (Cognitive): Information processing, numeracy, biases
    \item $\Psi_K$ (Informational): Transparency, signal quality, asymmetries
    \item $\Psi_E$ (Economic): Development, market thickness, resources
    \item $\Psi_T$ (Temporal): Time horizons, uncertainty, stability
    \item $\Psi_M$ (Market Scope): Geographic reach, trade access, openness
    \item $\Psi_F$ (Factor Flexibility): Adjustment capacity, mobility
\end{itemize}

\textbf{Empirical validation:} Across six canonical economic relationships (education$\rightarrow$earnings, management$\rightarrow$productivity, democracy$\rightarrow$growth, patience$\rightarrow$growth, information$\rightarrow$behavior, income$\rightarrow$health), the eight $\Psi$ dimensions explain 70.1\% of cross-country heterogeneity (Adjusted $R^2$ = 55.2\%). Different dimensions dominate for different effects: $\Psi_F$ (flexibility) matters most for management practices; $\Psi_K$ (information) matters most for education returns \citep{hanushek2012}; $\Psi_T$ (temporal stability) matters most for patience effects.

\textbf{Methodological distinction from LLM-based approaches:} This framework differs fundamentally from recent attempts to simulate human behavior using large language models (LLMs). \citet{peng2025} showed that ``digital twins''---LLMs conditioned on rich individual-level data---replicate only 50\% of classic behavioral economics experiments, often reproducing textbook predictions that humans violate. The present framework takes the opposite approach: instead of training on text \emph{about} experiments, we calibrate on the experimental \emph{data themselves}. Fifty years of behavioral economics research provides not mere illustration but \emph{calibration data} for the reference structure $C^*(\Psi)$. This yields functional knowledge---$Y = f(X, \Psi)$ with estimated parameters---rather than propositional knowledge (``effect X exists'').

\textbf{Self-improving research architecture:} The framework is not static. It incorporates two learning mechanisms that LLMs lack: (1)~\emph{Cumulative learning}---each new experiment improves $C^*(\Psi)$, whereas LLMs require complete retraining to incorporate new evidence; and (2)~\emph{Self-reinforcement learning}---when predictions deviate from observations ($\Delta = C - C^* \neq 0$), these deviations trigger a diagnostic protocol (Appendix~AK) that classifies the error type and translates it into structured model revision (Appendix~AL). This feedback loop---analogous to how AlphaZero improves through self-play and numerical weather prediction improves through forecast-error assimilation---transforms the framework from a model into a \emph{self-improving scientific infrastructure}. The traditional theory-test-revise cycle, which takes decades in economics, can be compressed to months. See Section~4.X for the full argument.

Welfare is measured through the FEPSDE framework: six dimensions (Financial, Emotional, Physical, Social, Digital, Ecological), three utility categories (Individual, Collective, Identity), four time horizons, and gains versus pains---yielding 144 components (detailed in Appendix~C) that capture human flourishing beyond GDP.

Novel predictions include: (1) treatment effects are context-dependent functions, not universal constants---$\tau(X \rightarrow Y) = f(\Psi)$; (2) information interventions work when $\Psi_K$ is \emph{low} (novel information) AND $\Psi_F$ is \emph{high} (can act on it); (3) trade wars have persistent effects through context hysteresis; (4) financial crises cause political crises with 5-10 year lags; (5) recovery time is convex in crisis depth.

The framework's core falsifiable claim: \emph{No economic effect is context-free.} If any intervention produces identical effects across all $\Psi$ configurations, the framework would be seriously challenged. This distinguishes our approach from theories that explain everything post hoc---we commit to context-dependence as a testable hypothesis.

\vspace{0.5em}
\noindent\textbf{Keywords:} Complementarity, Context, Metatheory, Welfare, FEPSDE, Institutional Economics, Behavioral Economics, Supermodularity, Treatment Effect Heterogeneity, Behavioral Calibration, Self-Reinforcement Learning

\vspace{0.5em}
\noindent\textbf{JEL Classification:} B41 (Economic Methodology), C62 (Existence and Stability), D01 (Microeconomic Behavior), D02 (Institutions), E02 (Institutions and the Macroeconomy), O43 (Institutions and Growth)
\end{abstract}

%%%%%%%%%%%%%%%%%%%%%%%%%%%%%%%%%%%%%%%%%%%%%%%%%%%%%%%%%%%%%%%%%%%%%%%
